\clearpage
\VerdierPhysicalPageAnchor{0085}%
\VerdierSemanticAnchor{ch1-3-1-4-homotopy-definition-continuation}
be two morphisms of $\operatorname{tr}_1(C)$. A \emph{homotopy} connecting
$(g,h,a)$ to $(g',h',a')$ is a triple $(b,b',\alpha)$, where
$b\colon g\rightsquigarrow g'$ and $b'\colon h\rightsquigarrow h'$ are
homotopies and where

\VerdierSemanticAnchor{ch1-3-1-4-second-order-condition}
\begin{equation*}
  \alpha\colon f'\star b+a'\rightsquigarrow b'\star f+a
\end{equation*}
is a second-order homotopy (2.5.12).

\VerdierSemanticAnchor{ch1-3-1-4-homotopic-morphisms}
Two morphisms of $\operatorname{tr}_1(C)$ will be said to be \emph{homotopic}
if there exists a homotopy connecting them.
\VerdierSemanticAnchor{ch1-3-1-4-equivalence-relation}
The relation ``$(g,h,a)$ and $(g',h',a')$ are homotopic'' is an equivalence
relation compatible with direct sum and composition of morphisms.
\VerdierSemanticAnchor{ch1-3-1-4-tr1-homotopy-category}
We may therefore introduce the category $\operatorname{Tr}_1(C)$, whose
objects are the objects of $\operatorname{tr}_1(C)$ and whose morphisms are the
morphisms of $\operatorname{tr}_1(C)$ up to homotopy.

\VerdierSemanticAnchor{ch1-3-1-5}
\noindent\textbf{Proposition 3.1.5.}\quad
\emph{Let
$D,D'\colon (X\xrightarrow{f}Y)\rightrightarrows
(X'\xrightarrow{f'}Y')$ be two morphisms of $\operatorname{tr}_1(C)$. If the
morphisms $D$ and $D'$ are homotopic in $\operatorname{tr}_1(C)$, then the
morphisms $c(D)$ and $c(D')$ are homotopic in $\operatorname{comp}(C)$.}

\VerdierSemanticAnchor{ch1-3-1-5-proof-setup}
Let $D=(g,h,a)$ and $D'=(g',h',a')$. Suppose that $D$ and $D'$ are homotopic
in $\operatorname{tr}_1(C)$. Thus there exist homotopies $b$ and $b'$ such
that:

\VerdierSemanticAnchor{ch1-3-1-5-eq1}
\begin{equation}
  \begin{aligned}
    g'-g&=bd+db,\\
    h'-h&=b'd+db'.
  \end{aligned}
  \tag{3.1.5.1}
\end{equation}

(To lighten the notation, in this calculation we have omitted the exponents
indicating the degree of the components.)
\VerdierSemanticAnchor{ch1-3-1-5-second-order-homotopy}
Moreover, there exists a second-order homotopy $\alpha$ such that:
\begin{equation*}
  \alpha d-d\alpha=b'f+a-(f'b+a').
\end{equation*}

\VerdierSemanticAnchor{ch1-3-1-5-cone-homotopy}
An immediate calculation then shows that
\begin{equation*}
  \begin{pmatrix}
    b' & \alpha\\
    0  & -b
  \end{pmatrix}
\end{equation*}
is a homotopy connecting $c(D)$ to $c(D')$.

\VerdierSemanticAnchor{ch1-3-1-6}
\noindent\textbf{Corollary 3.1.6.}\quad
\emph{The functors}
\begin{equation*}
  c\colon\operatorname{tr}_1(C)\longrightarrow\operatorname{comp}(C)
  \qquad\text{and}\qquad
  s\colon\operatorname{comp}(C)\longrightarrow\operatorname{tr}_1(C)
\end{equation*}
\emph{are compatible with the homotopies of $\operatorname{tr}_1(C)$ and
$\operatorname{comp}(C)$ and, upon passing to the quotient, define functors
still denoted}

\VerdierSemanticAnchor{ch1-3-1-6-descended-functors}
\begin{equation*}
  c\colon\operatorname{Tr}_1(C)\longrightarrow K(C)
  \qquad\text{and}\qquad
  s\colon K(C)\longrightarrow\operatorname{Tr}_1(C).
\end{equation*}
\VerdierSemanticAnchor{ch1-3-1-6-adjunction}
\emph{The functor $c\colon\operatorname{Tr}_1(C)\longrightarrow K(C)$ is left
adjoint to the functor
$s\colon K(C)\longrightarrow\operatorname{Tr}_1(C)$.}

\VerdierSemanticAnchor{ch1-3-1-6-proof}
This follows formally from (3.1.5) and (3.1.3).

\VerdierSemanticAnchor{ch1-3-1-7}
\noindent\textbf{3.1.7.}\quad
\VerdierSemanticAnchor{ch1-3-1-7-s-prime}
Denote by
$s'\colon\operatorname{comp}(C)\longrightarrow\operatorname{tr}_1(C)$ the
functor $X\longmapsto(X\longrightarrow 0)$.
\VerdierSemanticAnchor{ch1-3-1-7-comp-translation}
The functor
$cs'\colon\operatorname{comp}(C)\longrightarrow\operatorname{comp}(C)$ is
equal to the translation functor of $\operatorname{comp}(C)$.
\VerdierSemanticAnchor{ch1-3-1-7-descended-s-prime}
The functor $s'$ sends two homotopic morphisms of $\operatorname{comp}(C)$ to
homotopic morphisms of $\operatorname{tr}_1(C)$ and consequently defines a
functor, which we shall still denote by
$s'\colon K(C)\longrightarrow\operatorname{Tr}_1(C)$.
\VerdierSemanticAnchor{ch1-3-1-7-k-translation}
The functor $cs'\colon K(C)\longrightarrow K(C)$ is equal to the translation
functor of $K(C)$.
